\documentclass[11pt, letterpaper]{article}

% --- Packages ---
\usepackage[utf8]{inputenc}
\usepackage[T1]{fontenc}
\usepackage{lmodern}
\usepackage[margin=1in, headheight=14pt]{geometry}
\usepackage{amsmath, amssymb, amsthm, mathtools}
\usepackage{fancyhdr}
\usepackage{hyperref}

% --- Hyperref Setup ---
\hypersetup{
    colorlinks=true,
    linkcolor=black,
    urlcolor=blue,
    pdftitle={Project Euler Mathematical Notes}
}

% --- Theorem Environments ---
\theoremstyle{plain}
\newtheorem*{theorem}{Theorem}
\newtheorem*{lemma}{Lemma}
\newtheorem*{proposition}{Proposition}
\newtheorem*{corollary}{Corollary}

\theoremstyle{definition}
\newtheorem*{definition}{Definition}
\newtheorem*{observation}{Observation}

\theoremstyle{remark}
\newtheorem*{remark}{Remark}
\newtheorem*{complexity}{Complexity \& Bounds}

% --- Page Style ---
\pagestyle{fancy}
\fancyhf{}
\fancyhead[L]{\textsc{Project Euler Notes}}
\fancyhead[R]{\nouppercase{\leftmark}}
\fancyfoot[C]{\thepage}
\renewcommand{\headrulewidth}{0.4pt}

% --- Custom Problem Command ---
% Usage: \eulerproblem{<Number>}{<Title>}
% Example: \eulerproblem{1}{Multiples of 3 or 5}
\newcommand{\eulerproblem}[2]{%
    \clearpage
    \phantomsection
    \section*{Problem #1: #2}%
    \addcontentsline{toc}{section}{Problem #1: #2}%
    \markboth{Problem #1: #2}{Problem #1: #2}%
    \noindent\textbf{Archive Link:} \href{https://projecteuler.net/problem=#1}{\texttt{https://projecteuler.net/problem=#1}}
    \par\medskip\hrule\bigskip
}

% --- Common Macros ---
\newcommand{\N}{\mathbb{N}}
\newcommand{\Z}{\mathbb{Z}}
\newcommand{\Q}{\mathbb{Q}}
\newcommand{\R}{\mathbb{R}}
\newcommand{\C}{\mathbb{C}}
\DeclareMathOperator{\lcm}{lcm}

% --- Document Info ---
\title{\textbf{Project Euler: Mathematical Notes \& Reductions}}
\author{}
\date{\today}

\begin{document}

\maketitle
\tableofcontents

% ==========================================
% PROBLEM ENTRIES START HERE
% ==========================================

\eulerproblem{1}{Multiples of 3 or 5}

\subsection*{Core Idea}
Instead of iterating through all $n < N$, we use the Principle of Inclusion-Exclusion paired with closed-form arithmetic series evaluation to compute the sum in $\mathcal{O}(1)$ arithmetic operations.

\begin{lemma}
Let $S_k(N)$ denote the sum of positive multiples of $k$ strictly less than $N$. Setting $m = \lfloor \frac{N-1}{k} \rfloor$, we have:
\[
    S_k(N) = \sum_{i=1}^{m} ki = k \frac{m(m+1)}{2}.
\]
\end{lemma}

\subsection*{Derivation}
By inclusion-exclusion, the sum of all multiples of $3$ or $5$ strictly below $N = 1000$ is given by subtracting the intersection (multiples of $\lcm(3, 5) = 15$):
\[
    \text{Total} = S_3(1000) + S_5(1000) - S_{15}(1000).
\]


\eulerproblem{2}{Even Fibonacci Numbers}

\subsection*{Core Idea}
Every third Fibonacci number is even. We can derive a direct second-order linear recurrence for the subsequence of even Fibonacci numbers, eliminating the need to compute odd terms or check parity.

\begin{proposition}
Let $F_1 = 1, F_2 = 1$, and $F_n = F_{n-1} + F_{n-2}$. Define $E_k = F_{3k}$ as the $k$-th even Fibonacci number. Then $E_k$ satisfies the recurrence:
\[
    E_k = 4E_{k-1} + E_{k-2} \quad \text{for } k \ge 3,
\]
with seed values $E_1 = 2$ and $E_2 = 8$.
\end{proposition}

\begin{proof}
Expanding the standard recurrence relations down to indices divisible by $3$:
\begin{align*}
    F_{3k} &= F_{3k-1} + F_{3k-2} \\
           &= (F_{3k-2} + F_{3k-3}) + F_{3k-2} = 2F_{3k-2} + F_{3k-3} \\
           &= 2(F_{3k-3} + F_{3k-4}) + F_{3k-3} = 3F_{3k-3} + 2F_{3k-4} \\
           &= 3F_{3k-3} + F_{3k-4} + (F_{3k-5} + F_{3k-6}) \\
           &= 4F_{3k-3} + F_{3k-6} = 4E_{k-1} + E_{k-2}. \qedhere
\end{align*}
\end{proof}

\end{document}
